% © 2025 Ing. David Jaroš — CC BY-NC-ND 4.0
%
% This work is licensed under a Creative Commons Attribution-NonCommercial-NoDerivatives
% 4.0 International License (CC BY-NC-ND 4.0).
%
% License History: Earlier drafts (up to v0.3) were released under CC BY 4.0.
% From v0.4 onward, all material is released under CC BY-NC-ND 4.0 to protect
% the integrity of the theoretical work during ongoing academic development.
%
% See LICENSE.md for full license text.

\documentclass[11pt,a4paper]{article}
\usepackage{amsmath, amssymb, amsthm}
\usepackage{hyperref}
\usepackage[a4paper, margin=2.5cm]{geometry}
\usepackage{tcolorbox}
\usepackage{longtable}
\usepackage{booktabs}
\usepackage{slashed}

\newtheorem{theorem}{Theorem}[section]
\newtheorem{proposition}[theorem]{Proposition}
\newtheorem{lemma}[theorem]{Lemma}
\newtheorem{corollary}[theorem]{Corollary}
\theoremstyle{definition}
\newtheorem{definition}[theorem]{Definition}
\theoremstyle{remark}
\newtheorem{remark}[theorem]{Remark}

\title{\textbf{Fermion Sector Derivations}\\
\large Unified Biquaternion Theory — Canonical Appendices Series}
\author{Ing.\ David Jaroš}
\date{March 2026}

\begin{document}
\maketitle
\tableofcontents
\bigskip

\begin{tcolorbox}[colback=blue!4!white,colframe=blue!50!black,title=Purpose]
This appendix consolidates and strengthens the fermion sector derivations in UBT:
spinor representation from $\mathbb{C}\otimes\mathbb{H}$, the Dirac equation,
and the gauge coupling structure.

\medskip
\textbf{Primary sources}:
\texttt{canonical/qm\_emergence/step3\_dirac\_emergence.tex},
\texttt{canonical/qm\_emergence/step6\_spinorial\_subspace.tex},
\texttt{canonical/chirality/},
\texttt{canonical/interactions/qed.tex}

\medskip
\begin{tabular}{lll}
\hline
Derivation & Status & Section \\
\hline
Pauli matrices from $\mathbb{C}\otimes\mathbb{H}$ & \textbf{Proved} & \S\ref{sec:pauli} \\
Dirac gamma matrices & \textbf{Proved} & \S\ref{sec:gamma} \\
Spinor representation & \textbf{Proved} & \S\ref{sec:spinor} \\
Dirac equation & \textbf{Sketch} & \S\ref{sec:dirac} \\
Gauge coupling & \textbf{Proved} & \S\ref{sec:coupling} \\
Three generations & \textbf{Algebraic} & \S\ref{sec:generations} \\
\hline
\end{tabular}
\end{tcolorbox}

% -----------------------------------------------------------------
\section{Pauli Matrices from $\mathbb{C}\otimes\mathbb{H}$}\label{sec:pauli}
% -----------------------------------------------------------------

\begin{theorem}[Pauli matrices from quaternion generators]\label{thm:pauli}
The algebra isomorphism $\mathbb{C}\otimes\mathbb{H} \cong \mathrm{Mat}(2,\mathbb{C})$
maps the quaternion imaginary units $(I,J,K)$ to the Pauli matrices:
\begin{equation}\label{eq:pauli_map}
  I \;\mapsto\; i\sigma_1,
  \qquad
  J \;\mapsto\; i\sigma_2,
  \qquad
  K \;\mapsto\; i\sigma_3,
\end{equation}
where
\begin{equation}
  \sigma_1 = \begin{pmatrix}0&1\\1&0\end{pmatrix},\quad
  \sigma_2 = \begin{pmatrix}0&-i\\i&0\end{pmatrix},\quad
  \sigma_3 = \begin{pmatrix}1&0\\0&-1\end{pmatrix}.
\end{equation}
\end{theorem}

\begin{proof}
The standard isomorphism sends $1 \mapsto \mathbf{1}_2$,
$i \mapsto i\mathbf{1}_2$ (complex unit), $I \mapsto i\sigma_1$, etc.
The commutation relations $[I,J]=2K$, $[J,K]=2I$, $[K,I]=2J$ of $\mathbb{H}$
map exactly to $[i\sigma_a, i\sigma_b] = -[\sigma_a,\sigma_b]
= -2i\epsilon_{abc}\sigma_c = 2(i\epsilon_{abc})(i\sigma_c)/(i)
= 2\epsilon_{abc}(i\sigma_c)$, reproducing quaternion relations.
$\square$
\end{proof}

\noindent\textbf{Source}: \texttt{canonical/qm\_emergence/step3\_dirac\_emergence.tex \S2--\S3}.

% -----------------------------------------------------------------
\section{Dirac Gamma Matrices}\label{sec:gamma}
% -----------------------------------------------------------------

\begin{theorem}[Gamma matrices from $\mathbb{B}\otimes\mathbb{B}$]\label{thm:gamma}
The Dirac gamma matrices in the Weyl (chiral) representation emerge from
the tensor product $(\mathbb{C}\otimes\mathbb{H}) \otimes (\mathbb{C}\otimes\mathbb{H})$:
\begin{equation}\label{eq:gamma_weyl}
  \gamma^0 = \sigma_1 \otimes \mathbf{1}_2,
  \qquad
  \gamma^i = i\sigma_2 \otimes \sigma_i,\quad i=1,2,3.
\end{equation}
These satisfy the Clifford algebra:
\begin{equation}\label{eq:clifford}
  \{\gamma^\mu, \gamma^\nu\} = 2\eta^{\mu\nu}\mathbf{1}_4,
\end{equation}
with $\eta^{\mu\nu} = \mathrm{diag}(+1,-1,-1,-1)$.
\end{theorem}

\begin{proof}
Direct computation: $(\gamma^0)^2 = (\sigma_1)^2 \otimes \mathbf{1}_2 = \mathbf{1}_4$,
$(\gamma^i)^2 = -(\sigma_2)^2 \otimes (\sigma_i)^2 = -\mathbf{1}_4$ for $i=1,2,3$,
$\{\gamma^0,\gamma^i\} = \sigma_1(i\sigma_2)\otimes\sigma_i + (i\sigma_2)\sigma_1\otimes\sigma_i
= i(\sigma_1\sigma_2+\sigma_2\sigma_1)\otimes\sigma_i = 0$
(since $\{\sigma_a,\sigma_b\}=2\delta_{ab}$ gives $\sigma_1\sigma_2+\sigma_2\sigma_1=0$).
Clifford relation~\eqref{eq:clifford} is verified. $\square$
\end{proof}

\noindent\textbf{Source}: \texttt{canonical/qm\_emergence/step3\_dirac\_emergence.tex \S4}.

The chiral (Weyl) projectors are:
\begin{equation}
  P_L = \tfrac12(\mathbf{1}-\gamma^5), \qquad
  P_R = \tfrac12(\mathbf{1}+\gamma^5), \qquad
  \gamma^5 = i\gamma^0\gamma^1\gamma^2\gamma^3.
\end{equation}

% -----------------------------------------------------------------
\section{Spinor Representation}\label{sec:spinor}
% -----------------------------------------------------------------

\begin{definition}[Spinorial subspace]
The spinorial subspace of $\Theta$ is the projection onto the two-component
sub-algebra corresponding to each chirality sector:
\begin{equation}
  \xi := P_L\,\Theta_s \in \mathbb{C}^2
  \quad (\text{left-handed Weyl spinor}),
\end{equation}
where $\Theta_s$ is the spinorial sector of $\Theta$ identified via the
isomorphism $\mathrm{Mat}(2,\mathbb{C}) \cong \mathbb{C}\otimes\mathbb{H}$.
\end{definition}

\noindent\textbf{Source}: \texttt{canonical/qm\_emergence/step6\_spinorial\_subspace.tex}.

\begin{proposition}[Dirac spinor from $\Theta$]\label{prop:dirac_spinor}
A Dirac spinor $\Psi = \begin{pmatrix}\xi_L \\ \xi_R\end{pmatrix}$
arises as the $(2\oplus 2)$-decomposition of the spinorial component of $\Theta$
under the chirality projectors $P_{L,R}$.
\end{proposition}

% -----------------------------------------------------------------
\section{Dirac Equation}\label{sec:dirac}
% -----------------------------------------------------------------

\begin{tcolorbox}[colback=yellow!8!white,colframe=orange!60!black,title=Current status: Sketch]
The massive Dirac equation from $S[\Theta]$ is currently a sketch.
The structural derivation is complete; the identification of mass and the
full Euler-Lagrange variational argument require further work.
\end{tcolorbox}

The free Dirac operator is:
\begin{equation}\label{eq:Dirac_op}
  \mathcal{D} = i\gamma^\mu\nabla_\mu,
\end{equation}
which is uniquely determined (up to gauge equivalence) from the bundle structure
of $\Theta$ over spacetime.
(\texttt{research\_tracks/legacy\_theory\_variants/ubt/operators/dirac\_like\_operator.tex Thm.~2.1}).

The free Dirac equation $\mathcal{D}\Psi = m\Psi$ in UBT context arises from the
free-field limit of the master equation $\nabla^\dagger\nabla\Theta = 0$ upon
restricting to the spinorial sector and applying a first-order factorisation:
\begin{equation}\label{eq:dirac_factorisation}
  \nabla^\dagger\nabla = (\mathcal{D} - m)(\mathcal{D} + m) + O(R),
\end{equation}
where $O(R)$ denotes curvature corrections.

\begin{tcolorbox}[colback=yellow!8!white,colframe=orange!60!black,title=Open gap: mass identification]
\textbf{Gap D1}: The mass $m$ in~\eqref{eq:dirac_factorisation} must be derived
from $V(\Theta)$ or from topological winding (soliton interpretation).
A sketch is in \texttt{canonical/qm\_emergence/step3\_dirac\_emergence.tex \S6}.
\end{tcolorbox}

% -----------------------------------------------------------------
\section{Gauge Coupling Structure}\label{sec:coupling}
% -----------------------------------------------------------------

The gauged Dirac equation follows from replacing $\nabla_\mu \to D_\mu$
in~\eqref{eq:Dirac_op}:
\begin{equation}\label{eq:Dirac_gauged}
  (i\gamma^\mu D_\mu - m)\Psi = 0,
\end{equation}
with $D_\mu = \partial_\mu + \Gamma_\mu + igA_\mu + ig_s G^a_\mu T^a$.

\begin{proposition}[Gauge coupling from $D_\mu$ in $S[\Theta]$]
\label{prop:coupling}
The gauge-covariant derivative $D_\mu$ acting on the spinorial sector of $\Theta$
is exactly the coupling appearing in the gauged Dirac equation~\eqref{eq:Dirac_gauged}.
The electromagnetic, weak, and color couplings $e$, $g_W$, $g_s$ are the
eigenvalues of the respective gauge generators on the spinorial subspace.
\end{proposition}

\begin{proof}
The kinetic term $\mathrm{Tr}[(D_\mu\Theta)^\dagger D^\mu\Theta]$ restricted to
the spinorial sector gives $\bar{\Psi}\gamma^\mu D_\mu\Psi$.  The coupling
constants appear as the embedding eigenvalues of the internal symmetry generators
($Q$, $T^a$, $\lambda^a$) in the respective representations. $\square$
\end{proof}

\noindent\textbf{Source}: \texttt{canonical/interactions/qed.tex \S2--\S3}.

% -----------------------------------------------------------------
\section{Three Generations}\label{sec:generations}
% -----------------------------------------------------------------

UBT offers two complementary algebraic origins for the three fermion generations:

\paragraph{Origin 1: Quaternionic imaginary units $\{I,J,K\}$.}
The three non-scalar irreducible components of $\Theta \in \mathbb{H}$
correspond to three families:
$\Theta = \Theta_0 + \Theta_I\,I + \Theta_J\,J + \Theta_K\,K$,
where $\Theta_0$ (scalar) = Higgs sector,
$\Theta_I, \Theta_J, \Theta_K$ = three fermion generations.
(\texttt{research\_tracks/research/generation\_structure.tex Prop.~2})

\paragraph{Origin 2: $\psi$-winding modes.}
Independent $\psi$-winding modes carry the same $\mathrm{SU}(3)$ quantum numbers
(since $\mathrm{SU}(3)$ commutes with $\psi$) and ψ-parity forbids
odd$\leftrightarrow$even mixing, providing a structural separation of generations.
(\texttt{experiments/research\_tracks/three\_generations/st3\_complex\_time\_generations.tex \S3--\S4})

\begin{tcolorbox}[colback=yellow!8!white,colframe=orange!60!black,title=Open: mass spectrum]
\textbf{Open Hard Problem}: Reproducing the lepton mass ratios
$m_\mu/m_e \approx 207$ and $m_\tau/m_\mu \approx 16.8$ remains unsolved.
The KK mismatch theorem (ratio $\leq 2$ for any real modulus) proves that the
KK mechanism alone cannot reproduce the ratios.
See \texttt{DERIVATION\_INDEX.md \S fermions}.
\end{tcolorbox}

% -----------------------------------------------------------------
\section{Summary}\label{sec:summary_fermions}
% -----------------------------------------------------------------

\begin{center}
\begin{longtable}{llll}
\toprule
Result & Status & Layer & Source \\
\midrule
$\mathbb{C}\otimes\mathbb{H} \cong \mathrm{Mat}(2,\mathbb{C})$ & Proved & L0 &
  \texttt{step3\_dirac\_emergence.tex \S2} \\
Pauli matrices $\sigma^a$ from $\mathbb{H}$ & Proved & L0 &
  Thm.~\ref{thm:pauli} \\
Dirac $\gamma^\mu$ from $\mathbb{B}\otimes\mathbb{B}$ & Proved & L0 &
  Thm.~\ref{thm:gamma} \\
Clifford algebra $\{\gamma^\mu,\gamma^\nu\}=2\eta^{\mu\nu}$ & Proved & L0 &
  Thm.~\ref{thm:gamma} \\
Dirac operator $\mathcal{D}$ (unique) & Proved & L0 &
  \texttt{dirac\_like\_operator.tex} \\
Spinorial subspace of $\Theta$ & Defined & L0 &
  \texttt{step6\_spinorial\_subspace.tex} \\
Gauge coupling from $D_\mu$ & Proved & L0 &
  Prop.~\ref{prop:coupling} \\
Massive Dirac eq.\ from $S[\Theta]$ & Sketch & L1 &
  \S\ref{sec:dirac} \\
Mass identification (Gap D1) & Open & L1 & — \\
Three generations (algebraic) & Algebraic & L1 &
  \S\ref{sec:generations} \\
Mass ratios $m_\mu/m_e$, $m_\tau/m_\mu$ & Open Hard Problem & L2 &
  \texttt{DERIVATION\_INDEX.md} \\
\bottomrule
\end{longtable}
\end{center}

\end{document}
